% Imporditud: tools/import_eksperimendid.py
% Allikas: Eesti füüsikaolümpiaadi eksperimendi ülesannete
%   kogu (Rael Kalda, 2023), ülesanne Ü138, lahendus L138.
\setAuthor{EFO žürii}
\setRound{lõppvoor}
\setYear{2009}
\setNumber{G E1}
\setDifficulty{4}
\setTopic{staatika}
\setVahendid{ümmargune varras, nöör, dünamomeeter, koormis, statiiv}
\setPunktid{14}
\setAllikas{Eksperimendi kogu Ü138, lahendus lk 104}
\setLahendusseis{toores}

\prob{Koormise mass}
Määrata koormise mass. Hinnata mõõtemääramatust.
\textit{Märkus:} koormise kaal ületab dünamomeetri maksimaalse näidu.

\hint

\solu
Paneme klotsi nööri otsa rippuma nii, et nöör teeb pool ringi ümber pliiatsi ja

hoiame nööri vaba otsa paigal dünamomeetri abil; olgu lugem $F_1$. Nüüd kordame katset selle erinevusega, et nöör teeb täisringi ümber pliiatsi; olgu lugem $F_2$. Et mg/$F_1$ = $F_1$/$F_2$, siis m = F12/F2g.
\probend
